\documentclass[10pt,a4paper]{amsart}
\usepackage[margin=19mm]{geometry}
\usepackage{amsmath,amssymb,mathtools}
\usepackage[scr=rsfs,cal=euler]{mathalfa}
\usepackage{xfrac}
\usepackage[unicode,hidelinks,bookmarks=false]{hyperref}
\newcommand{\ie}{i.e.\ }
\newcommand{\cf}{cf.\ }
\newcommand{\resp}{respectively\ }
\newcommand{\Subsection}{\subsection}
\providecommand{\nobreakdash}{\nobreak\hbox{-}}
\setlength{\emergencystretch}{2em}
\multlinegap0pt
\usepackage{bm}
\usepackage{bbm}
\usepackage{dsfont}
\usepackage{xspace}
\usepackage{cancel}
\usepackage{mathtools}
\usepackage{amsthm}
\makeatletter
\@ifundefined{lemma}{\newtheorem{lemma}{Lemma}}{}
\@ifundefined{theorem}{\newtheorem{theorem}{Theorem}}{}
\makeatother
\title[Non-Definability of Reachability in B\"uchi Arithmetic for a]{Non-Definability of Reachability in B\"uchi Arithmetic for a Family of Generalized Collatz Maps}
\author{Madhav Dhiman, Rohan Pandey}
\date{Source: arXiv:2602.06066v2 (CC BY 4.0)}
\begin{document}
\maketitle

\noindent\emph{Source and licence.} Non-Definability of Reachability in B\"uchi Arithmetic for a Family of Generalized Collatz Maps. Version arXiv:2602.06066v2 (CC BY 4.0), distributed under the Creative Commons Attribution 4.0 International licence, as recorded in the arXiv licence metadata for this exact version and shown on the abstract page https://arxiv.org/abs/2602.06066v2. Full rights notice: licenses/arxiv-2602.06066-CC-BY-4.0.txt.\par\smallskip

\noindent\textit{Editorial scope.} Counted unit: the Theorem whose source label is \texttt{thm:main} in the article (source0.tex lines 360--364, source label \texttt{thm:main}), delivered together with the article's own complete original proof of it (source0.tex lines 365--397, one \texttt{proof} environment of 33 lines). No mathematics is written, repaired or re-derived: statement and proof are the article's own text line for line. Required context retained uncounted: lem:undefinable\_subset (lines 183--187); lem:power\_two\_path (lines 269--273); lem:non\_power\_two\_path (lines 284--288); lem:extraction (lines 337--345). Every definition, display and label of those lines is retained unchanged. Omitted from this package and not redistributed: 1--182, 188--268, 274--283, 289--336, 346--359, 398--567. Nothing inside a retained excerpt is omitted or altered: every retained body is delivered exactly as the article's lines are written, so no omission receipt of any kind accompanies this package. Citations: the retained excerpts carry no citation command at all, so no bibliography entry of the article is transcribed and no reference is redistributed. Reference adaptation: none was needed and none was made; every reference inside the delivered unit resolves inside it, so no unresolved reference or citation is produced. Printed slips kept literal: no printed word, symbol, hypothesis or number of the counted statement or of its proof was altered. Layout and class: the article's own class is \texttt{llncs} with packages including fontenc, amsmath, amssymb, hyperref; the wrapper is the lane's pinned \texttt{amsart} editorial wrapper at 10pt on A4, which restates the theorem-like environments the retained text needs and adds the article's own macro definitions that the retained text needs (none), with extra wrapper packages (bm, bbm, dsfont, xspace, cancel, mathtools). The delivered numbering is the unit's own. Assets: no retained span contains a figure environment, a graphics-inclusion command, a PDF-page inclusion, a table or a TikZ picture; no image file is copied into this package, the package declares no asset dependency and no image file is redistributed with it. Licence: version arXiv:2602.06066v2 of this article is distributed under the Creative Commons Attribution 4.0 International licence, as recorded in the arXiv licence metadata for this exact version and shown on the abstract page https://arxiv.org/abs/2602.06066v2. A full rights notice is delivered at licenses/arxiv-2602.06066-CC-BY-4.0.txt. Characters: every retained excerpt is pure ASCII, so the delivered engine, XeLaTeX with its default Latin Modern text font, reports no missing character.
\begin{lemma}\label{lem:undefinable_subset}
Let $q \ge 3$ be an odd integer. The set of powers of $2$, $P = \{2^k
\mid k \in \mathbb{N}\}$, is not first-order definable in
$\mathrm{BA}_q$.
\end{lemma}

\begin{lemma}[Powers of $2$]\label{lem:power_two_path}
Assume $q+d = 2^s$. Let $x = 2^k$ for some integer $k \ge 0$. Then
$R(x,1)$ holds, and every $z$ with $R(x,z)$ and $z \ne 1$ is even. In
particular $D(x,1)$ holds.
\end{lemma}

\begin{lemma}[Non-powers of $2$]\label{lem:non_power_two_path}
Assume $q+d = 2^s$. Let $x \in \mathbb{Z}_{>0}$ not be a power of $2$,
and suppose $R(x,1)$ holds. Then there is an odd integer $m > 1$ with
$D(x,m)$.
\end{lemma}

\begin{lemma}\label{lem:extraction}
Assume $R(x,z)$ is first-order definable in $\mathrm{BA}_q$. Then the
formula
\[
\psi(x) \;\equiv\; D(x,1) \land \forall u\,\bigl( (D(x,u) \land u \neq
1) \implies \mathrm{Even}(u) \bigr)
\]
is first-order definable in $\mathrm{BA}_q$.
\end{lemma}

\begin{theorem}\label{thm:main}
Let $q \ge 3$ and $d \ge 1$ be odd integers with $q+d$ a power of $2$.
Then the reachability relation $R(x,z)$ of the generalized Collatz map
$T_{q,d}$ is not first-order definable in $\mathrm{BA}_q$.
\end{theorem}

\begin{proof}
Assume for contradiction that $R$ is first-order definable in
$\mathrm{BA}_q$. By Lemma~\ref{lem:extraction}, the set
\[
S = \{x \in \mathbb{Z}_{>0} \mid \psi(x)\}
\]
is first-order definable in $\mathrm{BA}_q$. We show $S$ is exactly the
set of powers of $2$.

\medskip
\noindent\textbf{Claim (a): $\{2^k \mid k \in \mathbb{N}\} \subseteq S$.}
Let $x = 2^k$. By Lemma~\ref{lem:power_two_path}, $D(x,1)$ holds, and
every reachable value $z \ne 1$ is even. Any $u$ with $D(x,u)$ satisfies
$R(x,u)$, so $u$ is reachable from $x$; if $u \ne 1$ then $u$ is even.
The universal clause of $\psi(x)$ therefore holds, so $2^k \in S$.

\medskip
\noindent\textbf{Claim (b): $S \subseteq \{2^k \mid k \in \mathbb{N}\}$.}
Let $x \in S$. Since $\psi(x)$ holds, $D(x,1)$ holds, so $R(x,1)$ holds.
Suppose for contradiction that $x$ is not a power of $2$. By
Lemma~\ref{lem:non_power_two_path}, there is an odd integer $m > 1$ with
$D(x,m)$. Since $m \ne 1$, the universal clause of $\psi(x)$ forces
$\mathrm{Even}(m)$, contradicting that $m$ is odd. Hence $x$ is a power
of $2$.

\medskip
\noindent\textbf{Contradiction.}
By Claims (a) and (b), $S = \{2^k \mid k \in \mathbb{N}\}$. Definability
of $R$ makes $S$ first-order definable in $\mathrm{BA}_q$, hence
$q$-recognizable. By Lemma~\ref{lem:undefinable_subset}, the set of
powers of $2$ is not $q$-recognizable. This is a contradiction.
Therefore $R$ is not first-order definable in $\mathrm{BA}_q$.
\end{proof}

\end{document}
